\section{Existing Literature on Household Credit and Buy Now, Pay Later}
\label{ch:literature}

\subsection{Agent-Based Models of Household Credit}
\label{sec:lit-abm}

Distress at the level of a population is more than the sum of distress at the level of a household,
which is the case for modelling it from the bottom up. Cardaci \cite{cardaci2018inequality} builds
a hybrid agent-based stock-flow-consistent model in which heterogeneous households and banks
interact through a credit network, and finds that the link between inequality, household debt and
financial instability runs through the structure of that network rather than through aggregate
leverage, with consumption following an \emph{expenditure cascade} rule so that a banking crisis
emerges endogenously. D'Orazio and Giulioni \cite{dorazio2017micro} reach a similar conclusion from
a different direction, unsustainable debt emerging from local rules rather than being written in.
Both supply the methodological premise of this thesis, since the mechanism under investigation, a
single household holding obligations to several lenders none of whom can observe the others, exists
only on the individual balance sheet and cannot be expressed by a representative agent.

Madeira \cite{madeira2018chile}, at the Central Bank of Chile, models households that absorb
labour-income shocks, choose among loan contracts, and default when they can no longer finance
minimum consumption. That default condition is a cash-flow insolvency test rather than a leverage
threshold, and it is adopted here because the \ac{NCA} affordability assessment is itself a
residual-income test. Madeira also validates against observed Chilean default rates, demonstrating
that a household-level credit \ac{ABM} in a middle-income economy can be held to external
behavioural targets.

Hamill et al. \cite{hamill2023creditcard} model the United Kingdom credit-card market and find that
promotional strategies chosen by profit-maximising lenders raise borrower indebtedness. Their model
also reproduces a non-monotonic relationship between income and debt-to-income ratio, with
middle-income households carrying the highest values, which is used here as an independent
validation target.

\subsection{Empirical Evidence on Buy Now, Pay Later}
\label{sec:lit-bnpl}

Tracking 10.6 million United States consumers, deHaan et al. \cite{dehaan2024bnpl} show that
adopting \ac{BNPL} leads to immediate increases in overdraft charges, credit-card interest and late
fees, concentrated among liquidity-constrained households, with causally identified estimates
putting the rise in overdraft fees at nearly 9\%. Their conclusion, that \ac{BNPL} facilitates
overborrowing among vulnerable consumers, is the primary empirical rationale for testing these
mechanisms in South Africa.

Where deHaan et al. establish that \ac{BNPL} damages financial health, the \ac{CFPB}
\cite{cfpb2025bnpl} establishes the shape of the exposure. From roughly 145 million applications
originated between 2017 and 2022, concurrent borrowing is normal: across 2021 and 2022, 63\% of
\ac{BNPL} borrowers held simultaneous loans at some point and 32\% held them across different
firms. Those two figures are the empirical foundation of the debt-stacking mechanism, since holding
several obligations at once, particularly across providers who cannot see one another, converts
what looks at the level of a single transaction like a small interest-free deferral into an
aggregate servicing burden that no individual lender has assessed.

Distress concentrates in the \ac{LMI} population \cite{hayashi2025constraints, wang2026buynow}.
Ackert et al. \cite{ackert2025bnpl} explain the demand: in an incentivised experiment consumers
chose \ac{BNPL} over a credit-card loan for the same purchase and expected their social networks to
approve, from which the authors conclude that a payment scheme carrying social approval raises
indebtedness. That claim licenses both the want-driven borrowing path and the peer channel of this
model.

\subsection{Behavioural Foundations for the Decision Rules}
\label{sec:lit-behavioural}

Experimentally elicited present bias predicts credit-card borrowing \cite{meier2010present}, which
justifies a borrowing trigger answering to immediate shortfall rather than to a lifetime budget
constraint, and the same argument carries to repayment, present-biased borrowers systematically
failing to execute the paydown they planned \cite{Kuchler2021}. Keys and Wang \cite{Keys2019} find
that 29\% of United States credit-card accounts regularly pay at or near the contractual minimum,
and that at least 22\% of near-minimum payers shift their payments when the formula changes, which
is a clear signal of anchoring. Scheduled amortisation is therefore rejected as a repayment rule
here, because a substantial minority of borrowers pay the salient amount instead of the amount that
retires the debt, so obligations persist and compound. Together these results describe an agent
that borrows on shortfall, repays by salience, and varies in how strongly it does both.

\subsection{The South African Regulatory Setting}
\label{sec:lit-regulatory}

The \ac{NCA} \cite{sa_nca_2005} obliges registered credit providers to conduct a mandatory
affordability assessment, prescribed in Regulation 23A \cite{sa_ncr_affordability_2015}. That
regulation prescribes no \ac{DSTI} ceiling. It prescribes a residual-income test: discretionary
income is gross income less statutory deductions, less a table of minimum necessary living expenses
by income band, less every other committed obligation the consumer discloses or the bureau records,
and affordability is judged against what survives that subtraction. The \ac{NCA} also caps the cost
of credit by sub-sector, with maximum rates in force from 6 May 2016 running to the repurchase rate
plus 21\% for unsecured credit transactions \cite{sa_ncr_ratecaps_2016}.

The exemption \ac{BNPL} providers claim from that regime is set out in
Section~\ref{sec:background}. Two points from it bear on the model. TransUnion
\cite{transunion_cps_sa_2025} finds \ac{BNPL} named by 20\% of consumers intending to apply for
credit, behind only credit cards at 30\% and personal loans at 28\%, which is the only available
indication of South African demand. And the invisibility of \ac{BNPL} obligations is a documented
feature of the regulatory perimeter rather than a simplification adopted for tractability, so
bureau visibility enters the model as a switchable lever rather than a fixed assumption.

\subsection{Synthesis and Gap}
\label{sec:lit-gap}

Each of these literatures answers a different part of the question and none answers it whole, as
Table~\ref{tab:gap} sets out.

\begin{table}[H]
\centering
\caption{What each literature establishes, and what it leaves open.}
\label{tab:gap}
\begin{tabular}{L{3.4cm}L{5.2cm}L{5.4cm}}
\toprule
\textbf{Literature} & \textbf{Establishes} & \textbf{Leaves open} \\
\midrule
Empirical \ac{BNPL} \newline \cite{dehaan2024bnpl, cfpb2025bnpl} & That \ac{BNPL} harms financial
health, and that concurrent borrowing across providers is normal at scale & Observed in a United
States market where the product is at least partly visible to bureaus, and described rather than
explained \\
\addlinespace
Behavioural \newline \cite{meier2010present, Keys2019, Kuchler2021, ackert2025bnpl} & Why
individuals borrow on shortfall, repay by salience, and choose a socially approved payment scheme &
Stops at the individual; says nothing about how those choices aggregate \\
\addlinespace
Household credit \ac{ABM} \newline \cite{cardaci2018inequality, dorazio2017micro,
madeira2018chile, hamill2023creditcard} & How individual credit behaviour aggregates into
population-level instability, and that such models can be held to external targets & None includes
a lender class exempt from affordability regulation and invisible to its competitors \\
\addlinespace
South African regulatory \newline \cite{sa_ncr_affordability_2015, nortonrose_bnpl_sa,
transunion_cps_sa_2025} & That exactly such a two-lender-type perimeter exists in law, and that
\ac{BNPL} is absent from the credit report & No estimate of what that perimeter does to household
balance sheets, and none obtainable from administrative data \\
\bottomrule
\end{tabular}
\end{table}

The South African regulatory perimeter \cite{sa_ncr_affordability_2015, nortonrose_bnpl_sa, transunion_cps_sa_2025} produces a credit
market containing two lender types, one regulated and informed and the other neither, lending to
the same households.